\section*{Capítulo \ref{ch:g3:stages-of-life} --- Las etapas de la vida}

\begin{solution}{exo:g3:stages-of-life:1}
Bebé, niño, \omterm{def:g3:stages-of-life:stages}{adolescente}, \omterm{def:g3:stages-of-life:stages}{adulto} y \omterm{def:g3:stages-of-life:stages}{vejez}.
\end{solution}

\begin{solution}{exo:g3:stages-of-life:2}
(a) En la infancia, a partir de los seis años más o menos. (b) En la
etapa adulta. (c) En la etapa de bebé --- el primer año es el más
rápido de todos.
\end{solution}

\begin{solution}{exo:g3:stages-of-life:3}
El estirón del bebé y el del \omterm{def:g3:stages-of-life:stages}{adolescente}, con el crecimiento tranquilo
y constante de la infancia entre los dos.
\end{solution}

\begin{solution}{exo:g3:stages-of-life:4}
La línea se queda plana: la edad sigue avanzando por abajo, pero la
altura ya no sube.
\end{solution}

\begin{solution}{exo:g3:stages-of-life:5}
Entre los $6$ y los $8$: de unos $\qty{115}{cm}$ a unos
$\qty{127}{cm}$ --- unos $\qty{12}{cm}$ en los dos años. Solo en el
primer año: de unos $\qty{50}{cm}$ a unos $\qty{75}{cm}$ --- unos
$\qty{25}{cm}$, el doble que esos dos años de infancia juntos.
\end{solution}

\begin{solution}{exo:g3:stages-of-life:6}
Porque una comparación solo es justa si las medidas se hacen de la
misma manera. Zapatos un día y calcetines al siguiente, o un libro
torcido, crearían un ``crecimiento'' (o lo esconderían) que nunca
ocurrió.
\end{solution}

\begin{solution}{exo:g3:stages-of-life:7}
Bellota (\omterm{def:g2:from-seed-to-plant:seed}{semilla}), \omterm{prop:g2:from-seed-to-plant:cycle}{plántula}, roble joven, roble \omterm{def:g3:stages-of-life:stages}{adulto} que florece y
da bellotas, roble viejo --- y, con el tiempo, la muerte.
\end{solution}

\begin{solution}{exo:g3:stages-of-life:8}
Respuesta abierta. Con nueve años se está en la etapa de niño; el
estirón de la adolescencia está todavía por delante.
\end{solution}

\begin{solution}{exo:g3:stages-of-life:9}
Mosca (semanas), gato (un año más o menos), persona (casi veinte
años), roble (varias decenas de años). En los cuatro, el orden de las
etapas es el mismo y no se puede saltar ni invertir: juventud, edad
adulta, \omterm{def:g3:stages-of-life:stages}{vejez}.
\end{solution}

\begin{solution}{exo:g3:stages-of-life:10}
No. A los nueve está en el tramo tranquilo de la infancia --- unos
$\qty{6}{cm}$ al \emph{año}, que es medio centímetro al mes:
muchísimo menos de lo que se puede ver entre dos medidas separadas por
un mes. La curva sube con seguridad, pero despacio; solo tiene que
esperar unos meses entre puntos.
\end{solution}
